% ===== 原文 PDF pp.64--67（印刷 pp.63--66）=====
% tr-06 / chunk a：4.5 的结尾（p.64）、图 11、§4.6 起（p.65）、图 12、图 13、例 4.4--4.7
% 本文件 tag：4.23（p.65）、4.24 与 4.25（p.66）

% [边界] p.64 首行只有「(4.12).」，它是 p.63 末行「…where each of the 4 punctures is a
% pole of the Hitchin field as in equation」那句的收尾，按「开始者写完」归 tr-05；
% tr-06 自本页 Exercise 4.3 起，不重译上文。

% [原文 p.63]
\begin{exercisebox}{习题 4.3}
  用图 8 中理论所描述的同一套步骤，计算图 10 中理论的 Seiberg--Witten 曲线。
  最终结果将在下一节的例 4.4 中给出。
\end{exercisebox}

% [Tutor Note: 习题 4.3 是把 6d 视角倒推回 4d 的练习。图 8 是 T_2 理论：3 个正则穿刺点的
% 球面（m_1,m_2,m_3 分别在 z = 0, 1, ∞），φ_2 = f(z,m_i)/[z^2(z-1)^2]，对应 SU(2)^3
% 三分基本表示下的一个自由 N=2 半超多重态；图 10 比它多一个穿刺点、多一个复结构模 q。
% 做法与 §4.4 完全平行——对每个穿刺点写出 Hitchin 场的极点 (4.12)，用 rank-one 的
% Casimir φ_2 = (1/2)Tr(Φ_z^2)（见 (4.11)）定出 φ_2 的极点结构，即得 (4.25)。]

\begin{center}
  \includegraphics[width=0.85\linewidth]{fig11.png}
\end{center}
\figtitle{11}{由带四个基本味的 $\cN=2$ $\SU(2)$ SQCD 所描述的 SCFT。左图中的每个
$\SU(2)$ 味节点对应图 10 中的一个穿刺点；右图使该 SCFT 完整的 $\SO(8)$ 味对称性变得显然。}

% [Tutor Note: 图 11 的左右两图是同一理论的两种箭图表示：左图把 SO(8) 味对称性拆成
% 4 个 SU(2)（每个对应图 10 一个穿刺点/一个质量），右图把 4 个基本超多重态合成一个
% SO(8) 味节点。圆圈是规范节点、方框是味节点，这是箭图的惯用画法。]

Seiberg--Witten 曲线将对应于一个带 4 个味的 $\SU(2)$ 箭图的曲线，如图 11 所示。
我们为同一个 SCFT 给出两个箭图，其理由与图 10 中的复结构模 $q$ 有关：一般说来，
UV 曲线的复结构模的\textbf{不同尖点极限}（cusp limit）会从 Class S 构造给出同一个 SCFT
的不同箭图表示。\footnote{注意，用 6d $A_1$ 理论、UV 曲线取图 10 的那种 Class S
构造，只让 $\SO(8)$ 味对称性的 $\SU(2)^4$ 子群成为显然的。还存在一种替代的 Class S
构造：用 6d $D_4$ 理论紧化在一个带一个正则完全穿刺点和一个非正则穿刺点的球面上
（见 §4.8.2），它使完整的 $\SO(8)$ 对称性变得透明。}
图 11b 展示的是带 4 个基本味超多重态的 $\SU(2)$ 规范理论通常所用的箭图。
然而，黎曼球面上 $q$ 的不同取值可能给出黎曼曲面的不同分解（factorization）。
图 12a 展示了 $q\to 0$ 极限下球面的分解。如我们此前所见，球面的每个穿刺点代表
一个 $\SU(2)$ 超多重态，而把两个球面上的一对穿刺点连接起来的管代表一个 $\SU(2)$
矢量多重态；并且 $q$ 与复化规范耦合的关系为 $q=\ee^{2\pi\ii\tau}$（因此 $q\to 0$
对应该共形规范理论的弱耦合极限）。于是这个理论对应于图 11a 中的箭图。如果我们换而
把 $q$ 调到 $q\to 1$（使得原来的规范理论变为强耦合），就得到图 12b 中的分解，那是
另一种规范理论描述，与前一种是 S 对偶的。在箭图的层面上，这个对偶交换的是对应
$\SO(8)$ 整体对称性不同 $\SU(2)$ 子群的质量参数。

% [Tutor Note: 为什么 q -> 0 是弱耦合？因为 q = e^{2πiτ}，而 τ = θ/2π + 4πi/g^2，
% 故 q -> 0 等价于 Im τ -> +∞ 即 g -> 0。反过来 q -> 1 需要 τ -> 0（强耦合）；
% 此时按另一种方式切开球面得到的仍是同一理论、只是弱耦合的变量换了，这正是
% S 对偶在 Class S 语言中的几何实现。]

% [原文 p.64]
\begin{center}
  \includegraphics[width=0.9\linewidth]{fig12.png}
\end{center}
\figtitle{12}{带 4 个穿刺点的黎曼球面的分解。}

% [Tutor Note: 图 12 中两个球各带 2 个穿刺点、用一段管相连：按 Class S 的读法，
% 「管」= 一个 SU(2) 规范节点，「穿刺点」= 味。故 (a) 读出的是两组 SU(2) 味
% （即图 11a 的四个 SU(2) 方框）；(b) 中质量标签 m_1…m_4 的重新分组，对应
% S 对偶下 SO(8) 内不同 SU(2) 子群的交换。]

请注意，这个对偶从 6d 的角度看似乎有些显然，但从图 11b 的 4d 箭图的角度看却是
高度非平凡的。我们要说的是：当取带四个基本味的 $\SU(2)$ SQCD 的强耦合极限时，
该理论看起来又重新是同一个 SQCD 的弱耦合版本，只不过交换了一些质量。这一现象随后
可以推广到 Class S 中的其他理论：做法是研究带穿刺点的黎曼曲面 $\cC$ 的复结构模空间
中的尖点，它们会在所得的 4d 理论中产生对偶 [33]。

\SUB{4.6}{Class S[$A_n$] 中 $\cN=2$ SCFT 的例子}

下面我们给出 Class S 中若干选定的 $\cN=2$ SCFT 的例子，用以说明。

\begin{examplebox}{例 4.4：}
  作为第一个例子，我们考虑 4d 中（也许）最简单的 SCFT：带四个基本味的 $\SU(2)$
  SQCD。这是 Class S 理论的原型例子，也可以在其他综述与教科书中找到，例如
  [62,64]。一位勤奋的、已经完成习题 4.3 的读者可以直接跳到下一个例子。

  出发点是 6d 的 $\cN=(2,0)$ $\fsu(2)$ SCFT，它紧化在一个带四个穿刺点的球面上，
  如图 10 所示。我们给每个穿刺点配上一个质量参数 $m_1,\dots,m_4$。这意味着在每个
  穿刺点的局部，Hitchin 场形如
  \begin{equation}
    \Phi_z \sim \frac{1}{z}
    \begin{pmatrix} m_i & 0 \\ 0 & -m_i \end{pmatrix}
    + \text{正则项}.
    \tag{4.23}
  \end{equation}
  我们处于秩一这一事实让日子好过：把行列式 (4.4) 展开成 (4.9) 的形式，便得到
  % [原文 p.65]
  \begin{equation}
    x^2 = \phi_2(z)\,.
    \tag{4.24}
  \end{equation}
  在每个穿刺点 $z=0,1,q,\infty$ 处施加条件 (4.23)（为了研究最后一个穿刺点，
  应当换到球面的另一个坐标卡 $z\to 1/w$），就得到一个单参数族的曲线，
  \begin{equation}
    x^2 = \frac{1}{z(z-1)(z-q)}
    \left(\frac{q}{z}\,m_1^2 + \frac{q(q-1)}{z-1}\,m_2^2
    + \frac{z-q}{z-1}\,m_3^2 + z\,m_4^2 - u\right).
    \tag{4.25}
  \end{equation}
  这就是共形 $\SU(2)$ SQCD 的 Seiberg--Witten 曲线，其中 $u\in\C$ 是库仑分支的
  常用参数。
\end{examplebox}

% [Tutor Note: (4.23) 的留数矩阵必须无迹——这正是 su(2) 的条件；其本征值 ±m_i 就是
% 该穿刺点处的质量参数，即 SU(2) 味对称性 Cartan 的取值。秩一时
% φ_2 = (1/2)Tr(Φ_z^2)（见 (4.11)），故正则穿刺点附近 φ_2 只有二阶极点、系数恰为
% m_i^2（这正是 (4.13)）：用 (4.25) 验证，z -> 0 时
% x^2 ≈ [1/(z·(-1)·(-q))]·(q m_1^2 / z) = m_1^2/z^2。]
% [Tutor Note: (4.24) 就是 rank-one 的 SW 曲线：x -> -x 是超椭圆对合，曲线是
% z 球面的二次覆盖，φ_2 的极点阶数决定穿刺点类型。(4.25) 中 u 是唯一不被质量参数
% 固定的模——rank-one 时库仑分支复一维，u = <tr φ^2> 是唯一的库仑分支算子。]

% （排版说明）原文把图 13 排在本页（印刷 p.65）顶端、恰在例 4.4 的续文之前；
% 为不把一个 examplebox 从中间切开，译文把图 13 紧接在例 4.4 之后（仍在本页段内、
% 且在下一个图之前，符合 STYLE_GUIDE §6「图放在原文图附近」）。
\begin{center}
  \includegraphics[width=0.92\linewidth]{fig13.png}
\end{center}
\figtitle{13}{黎曼曲面与穿刺点形状的例子。}

% [Tutor Note: 图 13 三个子图说明「穿刺点形状」如何编码味对称性：(a) 上排两个一行
% 的杨图 = 极大（完全）穿刺点，携带完整 SU(N)，下排两个「一列 N-1 格 + 一列 1 格」
% = 极小穿刺点，只携带 U(1)；(b) 三个极大穿刺点的球面给出 T_N；(c) 带一个正则
% 穿刺点的环面（亏格 1），对应 N=4 SU(2) SYM。杨图的行长/列长就是该穿刺点处
% Hitchin 极点按 su(2) 嵌入分解后各块的大小。]

% [原文 p.66]
\begin{examplebox}{例 4.5：}
  第二个例子如图 13a 所示。这是用 Class S 描述写出的 $\SU(N)$ 共形 SQCD：
  黎曼曲面 $\cC$ 是一个带 4 个穿刺点的球面。上面两个穿刺点是极大穿刺点（即杨图
  只有一行），它们携带完整的 $\SU(N)$ 味对称性；下面两个穿刺点则是极小的（即杨图
  有一列含 $N-1$ 个格子、一列含 1 个格子），其味对称性只有 $\mathrm{U}(1)$。
\end{examplebox}

% （排版说明）脚注 43 在原文里落在「例 4.6」的蓝框之内。tcolorbox 内部的 \footnote
% 不会排到页面底部（脚注文字会丢），故此处用 \footnotemark 在框内打标记、紧接框后用
% \footnotetext 排出脚注文字；编号仍由 LaTeX 自动连续给出，与原文的条数、位置一致。
\begin{examplebox}{例 4.6：}
  更一般地，Class S 构造给出的许多 SCFT 并没有拉格朗日量描述。图 13b 中的
  $T_N$ 理论就是一例。黎曼曲面同样是一个球面，但这次我们放上三个极大穿刺点。
  这个相互作用 SCFT 没有拉格朗日量描述\footnotemark，但它具有显然的
  $\SU(N)\times\SU(N)\times\SU(N)$ 整体对称性。要拥有 $\cN=2$ 的共形拉格朗日量
  描述，一个必要条件是：黎曼曲面上存在某些我们可以调节的复结构模（它们对应复化
  规范耦合）。图 10 中带 4 个穿刺点的球面就是这种情况，而只有 3 个穿刺点的 $T_N$
  理论则不是。
\end{examplebox}
\footnotetext{我们的意思是：不存在能用 $\cN=2$ 语言写出的拉格朗日量。有可能存在
一些 $\cN=1$ 的拉格朗日量，它们沿流流向 4d 的一个相互作用 $\cN=2$ 理论 [133]。
超对称沿该流被明显破坏，但在红外恢复。这类流有很多例子，也被应用到（我们将在
§5.2 中引入的）$(G,G')$ 理论，例如参见 [134,202--211]。}

% [Tutor Note: 脚注 43 的要点：这里的「无拉格朗日量」只指「没有 N=2 语言可写的
% 拉格朗日量」，不等于「理论没有 UV 完成」——它可以由 N=1 理论沿 RG 流到达，
% 只是沿途超对称明显破缺、到 IR 才恢复；这是引用 [133] 那类工作的动机。]

% [Tutor Note: 例 4.6 中「可调复结构模 => 边缘耦合 => 有拉格朗日量」只是必要条件。
% 3 穿刺点球面没有复结构模（三点可用 SL(2,C) 全固定到 0,1,∞），拿不到复化规范耦合；
% 4 穿刺点时第四点的位置 q 才是模，故只有它有边缘耦合。]

\begin{examplebox}{例 4.7：}
  黎曼曲面 $\cC$ 一般可以有更高的亏格。例如，当 $\cC$ 是图 13c 中那样的单穿孔环面
  时，由 $A_1$ 型 $\cN=(2,0)$ 理论在该环面上（带一个正则极小穿刺点）实现的 4d
  理论，对应于 $\cN=4$ $\SU(2)$ SYM（其保持 $\cN=2$ 的质量形变理论也被称为
  $\cN=2^*$ 理论）。极点的结构由式 (4.12) 给出。

  此外，同一个 4d SCFT 可以有多种 Class S 构造，它们对应的 UV 曲线各不相同。再回到
  $\cN=4$ $\SU(2)$ SYM 这个例子：正如我们将在 §4.8.3 中更详细地看到的，还存在一种
  替代的 Class S 构造，即用 $A_2$ 型 $\cN=(2,0)$ 理论紧化在一个带一个扭曲正则穿刺点
  和一个扭曲非正则穿刺点的球面上（如图 15）。这第二种构造使 Witten 的反常 [212]
  对于 $\cN=4$ SYM 的 $\SU(2)$ 对称性变得显然 [56,213]。
\end{examplebox}

% [Tutor Note: 例 4.7 的 N=2* 是 N=4 SYM 加上保持 N=2 的超对称的质量项：N=4 矢量
% 多重态分解为 N=2 矢量多重态 + 伴随超多重态，给后者一个质量。之所以说「保持 N=2」，
% 是因为 N=4 的超对称不允许任意加质量项，只有加到 N=2 伴随超多重态上的这种项才
% 仍保持 N=2。]

% [Tutor Note: 单穿孔环面的模空间复一维（就是环面的 τ），对应 N=4 SYM 的共形流形；
% §4.8.3 中「一个扭曲正则 + 一个扭曲非正则」穿刺点的球面构造，则用另一组几何变量
% 重写同一理论。两种构造互补——这正是例 4.9 末尾到 p.76 首句要说的事。]

